% ============================================================
%  Section III — The Relational Coherence Theorem
%  "The Relational Coherence Theorem"
%  Dr. Joshua Adams | 2026
%
%  STATUS: Definitions complete. Propositions stubbed.
%          Full development in next pass.
% ============================================================

\section{\label{sec:rct}The Relational Coherence Theorem}

We now formalize the concepts developed in Sec.~\ref{sec:background}.
The Relational Coherence Theorem (\RCT) is a four-proposition
framework that unifies the relational ontology of
Sec.~\ref{sec:rqm} with the decoherence dynamics of
Sec.~\ref{sec:decoherence} through a single quantitative object:
the Affinity Tensor $\alpha(S,E)$.
We begin with the definitions that underpin all four propositions,
drawing a careful distinction between the full tensorial object and the
scalar summary statistic, a distinction that will be essential in
Sec.~\ref{sec:ura} when we parameterize the $\URA$ gate.

% ── III.A Formal Definitions ──────────────────────────────────────────────────

\subsection{\label{sec:definitions}Formal Definitions}

\begin{definition}[Relational Density Matrix]
\label{def:rho}
Let $S$ and $E$ be quantum systems with Hilbert spaces
$\mathcal{H}_S$ and $\mathcal{H}_E$ respectively, and let
$\rho_{SE} \in \mathcal{L}(\mathcal{H}_S \otimes \mathcal{H}_E)$
be the joint density matrix.
The \emph{relational density matrix of $S$ given $E$} is

\begin{equation}
\label{eq:rho-def}
\rhoSE \;\equiv\; \rho_{SE},
\end{equation}

understood as the state of $S$ defined only in the context of
its relational network with $E$.
The reduced density matrix $\rho_S = \mathrm{Tr}_E[\rhoSE]$
represents the locally accessible information about $S$
after $E$'s degrees of freedom are traced out.
\end{definition}

\begin{definition}[Affinity Tensor]
\label{def:alpha-tensor}
Let $\{\vert \pi_k \rangle\}$ be the pointer basis of $S$
selected by $H_{\mathrm{int}}$ via einselection
(Sec.~\ref{sec:decoherence}).
The \emph{Affinity Tensor} is the off-diagonal part of the
reduced density matrix in this basis:

\begin{equation}
\label{eq:alpha-tensor}
\alpha(S,E) \;\equiv\; \rho_S \;-\; \sum_k
  \langle \pi_k \vert \rho_S \vert \pi_k \rangle\,
  \vert \pi_k \rangle\langle \pi_k \vert,
\end{equation}

\noindent
i.e., $\alpha(S,E)$ retains only the coherences
$\langle \pi_i \vert \rho_S \vert \pi_j \rangle$ for $i \neq j$
and zeros out the diagonal populations.
The matrix $\alpha(S,E)$ is Hermitian, traceless, and vanishes
if and only if $\rho_S$ is diagonal in the pointer basis --- the
condition of complete decoherence.
\end{definition}

\begin{definition}[Scalar Affinity]
\label{def:alpha-scalar}
Let $d = \dim \mathcal{H}_S$.
The \emph{scalar affinity} is the normalized Frobenius norm of
the Affinity Tensor:

\begin{equation}
\label{eq:alpha-scalar}
\bar{\alpha}(S,E)
  \;\equiv\;
  \sqrt{\frac{d}{d-1}}\;
  \bigl\lVert \alpha(S,E) \bigr\rVert_F
  \;=\;
  \sqrt{\frac{d}{d-1}\sum_{i \neq j}
    \bigl\lvert
      \langle \pi_i \vert \rho_S \vert \pi_j \rangle
    \bigr\rvert^2\;},
\end{equation}

\noindent
where $\lVert \cdot \rVert_F$ denotes the Frobenius norm.
The prefactor $\sqrt{d/(d-1)}$ normalizes $\bar{\alpha}$ to the
unit interval: $\bar{\alpha} \in [0, 1]$.
\end{definition}

\begin{remark}[Boundary values of $\bar{\alpha}$]
\label{rem:boundary}
For $d = 2$ (qubit), Eq.~\eqref{eq:alpha-scalar} reduces to
$\bar{\alpha} = 2\lvert \langle \pi_0 \vert \rho_S \vert \pi_1 \rangle \rvert$.
A pure qubit state $\vert \psi \rangle = \cos\theta\,\vert\pi_0\rangle
+ e^{i\phi}\sin\theta\,\vert\pi_1\rangle$ yields
$\bar{\alpha} = \lvert \sin 2\theta \rvert$, which attains its
maximum $\bar{\alpha} = 1$ at the equal-weight superposition
$\theta = \pi/4$ and its minimum $\bar{\alpha} = 0$ at either
pointer state.
A fully decohered mixture $\rho_S = p\,\vert\pi_0\rangle\langle\pi_0\vert
+ (1-p)\,\vert\pi_1\rangle\langle\pi_1\vert$ gives $\bar{\alpha} = 0$
for all $p$, consistent with the classical limit.
\end{remark}

\begin{remark}[Tensor vs.\ scalar]
\label{rem:tensor-vs-scalar}
The Affinity Tensor $\alpha(S,E)$ and the scalar affinity
$\bar{\alpha}(S,E)$ serve distinct roles throughout this paper.
$\alpha(S,E)$ is a matrix that encodes the full coherence
structure of $S$ --- the magnitude and phase of every off-diagonal
element --- and appears in the proof of Proposition~III
(Sec.~\ref{sec:prop3}) and in the parameterization of $\URA$
(Sec.~\ref{sec:ura}).
$\bar{\alpha}(S,E)$ is the single scalar that collapses this
structure to a coherence magnitude; it is the quantity that
transitions continuously from $1$ (maximally coherent, fully
quantum-relational) to $0$ (fully decohered, classically definite)
and that appears in the statement of all four propositions.
\end{remark}

% ── III.B–E  Propositions (full development: next pass) ──────────────────────

\subsection{\label{sec:prop1}Proposition~I: Property Indefiniteness}

% TODO: Full development in next pass.
% Claim: No quantum property of $S$ exists independently of a
% relational context defined by $\Gamma_S$.
% Proof strategy: cite Bell + Rovelli; show \bar{alpha} = 1
% is incompatible with a definite pointer-state value.

\subsection{\label{sec:prop2}Proposition~II: Relational Determination}

% TODO: Full development in next pass.
% Claim: Definite properties of $S$ emerge through the relational
% network \Gamma_S via decoherence-driven einselection.
% Proof strategy: Lindblad evolution + pointer-state stability.

\subsection{\label{sec:prop3}Proposition~III:
            Coherence--Affinity Correspondence}

% TODO: Full development in next pass.
% Claim: \bar{alpha} = 1 <=> maximal entanglement of rho_{SE};
%        \bar{alpha} = 0 <=> complete decoherence into pointer basis.
% Proof strategy: direct computation from Defs.~\ref{def:alpha-tensor}
% and \ref{def:alpha-scalar}; continuity argument for intermediate values.

\subsection{\label{sec:prop4}Proposition~IV:
            Conservation of Relational Potential}

% TODO: Full development in next pass.
% Claim: In a closed system S+E, total quantum information is conserved;
% decoherence of S transfers coherence to S-E correlations, not destroys it.
% Proof strategy: unitary evolution of rho_{SE}; purity argument;
% relate to quantum mutual information I(S:E).
